% Esquisse d'un programme (cote 119, copy 1): archive cote 119, site page 26 (PDF page 27), typescript page 24.
% typescript. Draft compared: EsquisseFr.pdf (text layer):867-942.
% AI-transcribed by Claude Opus 5.5 (claude-opus-5-5, adaptive thinking, effort high), Message Batches API, batch esq1 (msgbatch_01LNsogFyF3y6KSp9UbtJtBk), 2026-11-12; not yet checked by eye.
\scanpage{119/26}
\mspage{24}

pour l'icosaèdre, la caractéristique 2 pour l'octaèdre), on lit,
sur le polyèdre régulier géométrique sur le corps fini concerné
($\mathbb{F}_4$ et $\mathbb{F}_5$ pour l'icosaèdre, $\mathbb{F}_2$ pour l'octaèdre) une descrip-%
tion particulièrement élégante (et inattendue) de la combina-%
toire du polyèdre. Il m'a semblé même entrevoir là un principe
d'une grande généralité, que j'ai crû retrouver notamment dans
une réflexion ultérieure sur la combinatoire du système des 27
droites d'une surface cubique, et ses relations avec le système
de racines $E_7$. Qu'un tel principe existe bel et bien et qu'on
réunisse même à le dégager de son manteau de brumes, ou qu'il
recule au fur et à mesure où on le poursuit et qu'il finisse par
s'évanouir comme une Fata Morgana, j'y trouve pour ma part une
force de motivation, une fascination peu communes, comme celle du
rêve peut-être. Nul doute que de suivre un tel appel de l'infor-%
mulé, de l'informe qui cherche forme, d'un entrevu
élusif qui semble prendre plaisir à la fois à se dérober et à se
manifester - ne peut que mener loin, alors que nul ne pourrait
prédire, où ...

Pourtant, pris par d'autres intérêts et tâches, je n'ai
pas jusqu'à présent suivi cet appel, ni rencontré personne d'autre
qui ait voulu l'entendre, et encore moins le suivre. Mis à part
quelques digressions vers d'autres types de structures géométrico-
combinatoires, mon travail ici encore s'est borné à un premier tra-%
vail de dégrossissage et d'intendance, sur lequel il est inutile
de m'étendre plus ici\hand{$^{(5)}$}. Le seul point qui peut-être mérite encore
mention, est l'existence et l'unicité de l'hyperquadrique circons-%
crite à un n-polyèdre régulier donné, dont l'équation peut s'expli-%
citer par des formules simples en termes des paramètres fondamentaux
$\alpha_i$ (*). Le cas qui m'intéresse le plus est celui où $n = 2$, et le

\rule{5cm}{0.4pt}

(*) Un résultat analogue vaut pour les pseudo-polyèdres\unsure{.|:} Il semblerait
que les ``caractéristiques exceptionnelles'' dont il a été question plus
haut, pour les spécialisations d'un polyèdre donné, sont celles pour
lesquelles l'hyperquadratique circonscrite est, soit dégénérée, soit
tangente à l'hyperplan à l'infini.

\margin{26}
